\section{The CLDB System}
\label{sec:method}

\subsection{Overview}
CLDB is an online index-selection agent. Queries arrive continuously; each query
is parsed and encoded, the agent selects an action under a safety mask, the
action is validated and executed, the resulting latency is measured, and the
transition is written to a prioritized replay buffer. A background scheduler
performs gradient updates, and a drift detector triggers elastic weight
consolidation when the workload changes. Figure~\ref{fig:arch} summarizes the
pipeline. The design goal is that the \emph{critical path} (inference and
execution) never blocks on learning.

\begin{figure}[t]
\centering
\begin{tikzpicture}[
  font=\footnotesize, node distance=6mm and 7mm,
  box/.style={draw, rounded corners, align=center, inner sep=4pt, minimum height=7mm},
  arr/.style={-{Latex[length=2mm]}, thick}, bg/.style={draw, dashed, rounded corners, inner sep=6pt}
]
\node[box] (q) {Query $q_t$};
\node[box, right=of q] (p) {Parser\\[-2pt]\& Features};
\node[box, right=of p] (e) {Embedder};
\node[box, right=of e] (st) {State $s_t$};
\node[box, below=of st] (agent) {DQN\\[-2pt]$\pi_\theta$ (+ mask)};
\node[box, left=of agent] (safe) {Safety\\[-2pt]check};
\node[box, left=of safe] (exec) {Executor\\[-2pt](PG)};
\node[box, right=of agent] (buf) {Prioritized\\[-2pt]replay $\mathcal{D}$};
\node[box, below=of buf] (train) {Background\\[-2pt]DQN update};
\node[box, below=of train] (ewc) {Drift-triggered\\[-2pt]EWC};
\draw[arr] (q)--(p); \draw[arr] (p)--(e); \draw[arr] (e)--(st);
\draw[arr] (st)--(agent); \draw[arr] (agent)--(safe);
\draw[arr] (safe)--(exec); \draw[arr] (agent)--(buf);
\draw[arr] (buf)--(train); \draw[arr] (train)--(ewc);
\draw[arr, dashed] (ewc.west) -- ++(-1,0) |- (agent.south);
\begin{scope}[on background layer]
\node[bg, fit=(p)(e)(st)(agent)(buf)(train)(ewc), label={[font=\scriptsize\itshape]above:non-blocking learning path}] {};
\end{scope}
\end{tikzpicture}
\caption{The CLDB pipeline. Inference and execution form the critical path; the
replay buffer, background DQN updates and drift-triggered EWC run off the
critical path.}
\label{fig:arch}
\end{figure}

\subsection{State construction}
The state \eqref{eq:state} combines four heterogeneous signals. The
\emph{structural} features \(\phi(q)\) are cheap to compute and generalize across
table names via a fixed-size hashing trick. The \emph{semantic} embedding
\(e(q)\) is produced by a sentence encoder (with a deterministic hashing
embedder as an offline fallback, see Section~\ref{sec:experiments}); it lets the
agent generalize between textually similar queries. The \emph{workload context}
\(c_t\) summarizes recent activity. The \emph{index-state} \(\iota(I_t)\)
\eqref{eq:indexstate} closes the loop between decisions: it tells the agent
whether the current predicate is already covered and how much budget remains, so
that the value of \textsc{create} and \textsc{drop} can be assessed.

\subsection{Learning: off-policy DQN}
We approximate \(Q^\star\) with a network \(Q_\theta\) and a frozen target
network \(Q_{\theta^-}\). Transitions \((s_t,a_t,r_t,s_{t+1},d_t)\) are drawn
from the replay buffer \(\mathcal{D}\) where \(d_t\in\{0,1\}\) is the terminal
flag. The temporal-difference target and loss are
\begin{align}
  y_t &= r_t + \gamma\,(1-d_t)\,\max_{a'} Q_{\theta^-}(s_{t+1},a'), \label{eq:target}\\
  \loss_{\mathrm{DQN}}(\theta) &= \E_{(s,a,r,s',d)\sim\mathcal{D}}
     \Big[\, w \cdot \rho_\delta\big(y - Q_\theta(s,a)\big) \Big], \label{eq:dqnloss}
\end{align}
where \(\rho_\delta\) is the Huber loss and \(w\) is the importance-sampling
weight of \eqref{eq:is}. The target parameters are copied from \(\theta\)
every \(C\) updates. Gradients are clipped to norm \(10\).

\subsection{Database-aware prioritized replay}
A uniform buffer spends most of its capacity on the (many) common queries. We
instead prioritize experiences by their \emph{impact on the database}:
\begin{equation}
  p_i \propto \big(|\delta_i|+\epsilon\big)^{\alpha},
  \qquad
  \text{initial priority } \; \tilde p_i = \big(|r_i|+\epsilon\big)\cdot
      \tfrac{1}{\sqrt{f(a_i)}},
  \label{eq:priority}
\end{equation}
where \(\delta_i\) is the absolute TD error, \(f(a_i)\) is the running frequency
of action \(a_i\) in the buffer (so \emph{rare} decisions, e.g.\ analytical
\textsc{drop}s, are surfaced), and \(\epsilon\) is a small constant. Sampling uses
probability \(p_i\); the loss is corrected by the importance-sampling weight
\begin{equation}
  w_i = \frac{\big(N\,p_i\big)^{-\beta}}{\max_j (N\,p_j)^{-\beta}},
  \label{eq:is}
\end{equation}
with \(N=|\mathcal{D}|\). After each update, priorities are refreshed with the
observed TD error, \(\tilde p_i \leftarrow (|\delta_i|+\epsilon)\,/\sqrt{f(a_i)}\).
This combines the anti-forgetting effect of rehearsal with the sample efficiency
of prioritized replay, and adds a database-specific notion of importance (latency
impact and rarity) that generic replay lacks.

\subsection{Safety-constrained exploration}
Exploration uses \(\epsilon\)-greedy selection \emph{over the masked action set}
\(\mu(s_t)\). The mask forbids \textsc{create} without predicate columns (an index
on nothing is meaningless), \textsc{create} beyond the budget, and \textsc{drop}
of a non-existent index. Because the mask is applied before both exploration and
exploitation, the agent never proposes an infeasible or unsafe schema change,
which also removes a large amount of uninformative exploration. A post-hoc safety
layer additionally rejects low-confidence actions and can roll back a change whose
measured latency regresses.

\subsection{Drift-triggered elastic weight consolidation}
When the workload shifts, the policy should adapt---but not at the cost of the
phases already learned. EWC \citep{kirkpatrick2017ewc} adds a quadratic penalty
anchoring parameters in proportion to their importance for previous tasks,
\begin{equation}
  \Omega(\theta) = \frac{\lambda}{2}\sum_i F_i\,\big(\theta_i - \theta^\star_i\big)^2,
  \label{eq:ewc}
\end{equation}
where \(\theta^\star\) is the parameter snapshot at the end of the consolidated
phase and \(F\) is the Fisher information matrix (Section~\ref{sec:theory}). The
key design choice is \emph{when} to consolidate. Fixed task boundaries are
unknown online; instead we consolidate only when the drift detector
\citep{gama2014drift}, which tracks the cosine distance between the centroid of
a sliding window of query embeddings and a baseline centroid, fires. This
\emph{drift-triggered} rule performs consolidation \(O(m)\) times for \(m\) drift
events rather than \(O(T)\) times for \(T\) phases, without missing a shift
(Lemma~\ref{lem:consolidation}). Algorithm~\ref{alg:online} and
Algorithm~\ref{alg:ewc} summarize the online loop and consolidation.

\begin{algorithm}[t]
\caption{CLDB online step (one query)}
\label{alg:online}
\begin{algorithmic}[1]
\Require query $q_t$, policy $Q_\theta$, buffer $\mathcal{D}$, mask $\mu$
\State $s_t \gets \mathrm{build\_state}(q_t,\, I_t)$ \Comment{Eq.~\eqref{eq:state}}
\State finalize previous transition with successor $s_t$; push to $\mathcal{D}$
\State $a_t \gets \epsilon\text{-greedy}\big(Q_\theta(s_t),\,\mu(s_t)\big)$
\State validate $a_t$ with the safety layer; if rejected, $a_t\gets\textsc{keep}$
\State execute $a_t$; measure latency; compute $r_t$ \Comment{Eq.~\eqref{eq:reward}}
\State store $(s_t,a_t,r_t,s_{t+1})\to\mathcal{D}$
\State notify drift detector with $e(q_t)$; if drift, call \textsc{Consolidate}
\end{algorithmic}
\end{algorithm}

\begin{algorithm}[t]
\caption{Drift-triggered consolidation}
\label{alg:ewc}
\begin{algorithmic}[1]
\Require current states $\mathcal{B}$, phase id $p$, policy $Q_\theta$
\For{parameters $i$} $F_i \gets 0$ \EndFor
\For{batch $b \subset \mathcal{B}$}
  \State $\ell \gets \frac{1}{|b|}\sum_{s\in b}\sum_{a} Q_\theta(s,a)$
  \State $F_i \gets F_i + (\partial_{\theta_i}\ell)^2$
\EndFor
\State $F_i \gets F_i / \#\text{batches}$; \;
$\theta^\star_i \gets \theta_i$ \Comment{snapshot}
\State register $(p, F, \theta^\star)$ for EWC penalty \eqref{eq:ewc}
\end{algorithmic}
\end{algorithm}

\subsection{Non-blocking online operation}
Inference is a single forward pass (milliseconds). Training runs in a background
thread that wakes periodically and performs DQN updates from the replay buffer;
drift detection and EWC consolidation also run off the critical path. Query
execution therefore never blocks on learning, which is the operational
requirement that motivated the design. The pipeline exposes structured metrics
(replay size, TD error, EWC penalty, drift score).
